% !PRINTED 631
\weberpage{631}
%<<CONT
property that every root of an irreducible metacyclic equation of degree $n$ in $\Omega$ is contained in this expression, and that conversely every such expression also satisfies an irreducible equation of degree $n$ in $\Omega$.

The foundation for this investigation is furnished by the resolvents of \gesperrt{Lagrange}.

Now let $n$ be a prime, let $x_0, x_1, x_2 \ldots x_{n-1}$ be a system of independent variables, and let the notation be so chosen that $x_n = x_0$, $x_{n+1} = x_1 \ldots$ Finally, let $\varepsilon$ be an $n$th root of unity. We then set
\weq{1}{
(\varepsilon, x) = x_0 + \varepsilon x_1 + \varepsilon^2 x_2 + \cdots + \varepsilon^{n-1} x_{n-1} = \sum_{0, n-1}^{h} \varepsilon^h x_h.
}
These expressions we have already come to know in the fourteenth section under the name of the \gesperrt{resolvents of Lagrange}, and we have there (\S.\,164) investigated their behaviour with respect to the \gesperrt{cyclic permutations}.

It is now a matter of determining the influence of linear permutations, of the kind considered in \S.\,180, on these functions.

We first send ahead a lemma, which we really need only reproduce from what has gone before. The \gesperrt{imaginary} $n$th roots of unity $\varepsilon$ are the roots of an equation $X = 0$, if
\weq{2}{
X = \xi^{n-1} + \xi^{n-2} + \cdots + \xi + 1
}
is set, and $\xi$ denotes a variable.

The function $X$ is, as we have seen earlier (\S.\,134), irreducible in the field $R$ of the rational numbers.

If we take as basis any field $\Omega$ in which $X$ is irreducible, and derive from it the field $\Omega(\varepsilon)$, which consists of all rational functions of $\varepsilon$ with coefficients in $\Omega$, then it follows by \S.\,142 that every quantity of this field can be put in one and only one way into the form
\weq{3}{
c_0 + c_1 \varepsilon + c_2 \varepsilon^2 + \cdots + c_{n-2} \varepsilon^{n-2}
}
in which the $c$ are quantities in $\Omega$.

For brevity we shall call this form of the quantities in $\Omega(\varepsilon)$ the \gesperrt{normal form}, and shall accordingly state the theorem:

1. \gesperrt{If $\Omega$ is a field in which the cyclotomic equation $X = 0$ is irreducible, then every}
%CONT>>
